%\begin{itemize}\item Python being a high level language with a rich ecosystem of libraries for scientific computing and data processing, it is a natural choice for this project. Most of the routines are imlimented as python scripts. IPython notebooks are used for the development and testing of the routines.    \item PDAL is a C++ library for reading and writing point cloud data. It is used for the pre-processing of point cloud data. various filter including ISMR filter is used for the pre-processing of point cloud data. \item Laspy is a python library for reading and writing LAS files. It is used for the pre-processing of point cloud data.  \item GDAL, Rasterio and geopandas is used for the processing of raster data, managing vector data and for the reprojection of data. \item Numpy, Scipy and sklearn is used for the numerical computation and for the machine learning algorithms like knn, k-d tree, etc. \item Keras is a deep learning API used for the development of deep learning models. It is used for the training of deep learning models. \item Torch is a deep learning framework used for the development of deep learning models. It is used for the testing the Transformer based models. \item Redis is a in-memory data structure store, used as a distributed cache and message broker. It is used to orchester loads in simultaneous workers in a networked distributed system. \item Qgis is a free and open source geographic information system. It is used for the visualization and analysis of geographic data. \item Paraview is a open source data analysis and visualization application. It is used for the visualization of 3D data. \end{itemize}



\begin{itemize}
    \item \textbf{Python} is the primary language, with IPython notebooks used for development and testing.
    \item \textbf{PDAL} handles point cloud pre-processing, including SMRF and other filters.
    \item \textbf{Laspy} is used for reading and writing LAS files.
    \item \textbf{GDAL, Rasterio \& GeoPandas} manage raster processing, vector data, and reprojection.
    \item \textbf{NumPy, SciPy \& Scikit-learn} provide numerical computation and ML algorithms (KNN, k-d tree, etc.).
    \item \textbf{Keras \& PyTorch} are used for deep learning — Keras for CNN model training and PyTorch for Transformer-based model evaluation.
    \item \textbf{Redis} serves as a distributed cache and message broker to orchestrate load across parallel workers in a networked system.
    \item \textbf{QGIS} is used for visualization and analysis of geographic data.
    \item \textbf{ParaView} handles visualization of 3D point cloud data.
\end{itemize}


